\documentclass[varwidth=16cm, border=2mm]{standalone}
\usepackage{amsmath}\usepackage{amssymb}\usepackage{ctex}
\usepackage{tasks}\usepackage{xcolor}\usepackage{graphicx}
\settasks{label=\Alph*., label-width=1.5em, item-indent=2em}
\begin{document}
\textbf{[0.6]} 已知椭圆$\frac{x^2}{9} + y^2 = 1$，过左焦点$F$作倾斜角为$\frac{\pi}{6}$的直线交椭圆于$A、B$两点，则弦$AB$的长为_________.

\vspace{0.2cm}\par\noindent\textcolor{red}{\textbf{【答案】} 2}
\par\noindent\textcolor{blue}{\textbf{【解析】} 1. **确定椭圆参数和焦点坐标**：\n    椭圆方程为$\frac{x^2}{9} + y^2 = 1$，即$\frac{x^2}{3^2} + \frac{y^2}{1^2} = 1$。\n    因此，$a^2 = 9 \Rightarrow a = 3$， $b^2 = 1 \Rightarrow b = 1$。\n    根据椭圆的性质，$c^2 = a^2 - b^2 = 9 - 1 = 8$，所以$c = \sqrt{8} = 2\sqrt{2}$。\n    左焦点$F$的坐标为$(-c, 0)$，即$F(-2\sqrt{2}, 0)$。\n\n2. **确定直线的方程**：\n    直线过左焦点$F(-2\sqrt{2}, 0)$，倾斜角为$\frac{\pi}{6}$。\n    直线的斜率$k = \tan(\frac{\pi}{6}) = \frac{\sqrt{3}}{3}$。\n    根据点斜式，直线的方程为$y - 0 = \frac{\sqrt{3}}{3}(x - (-2\sqrt{2}))$，即$y = \frac{\sqrt{3}}{3}x + \frac{2\sqrt{6}}{3}$。\n\n3. **联立直线与椭圆方程**：\n    将直线方程$y = \frac{\sqrt{3}}{3}x + \frac{2\sqrt{6}}{3}$代入椭圆方程$\frac{x^2}{9} + y^2 = 1$：\n    $\frac{x^2}{9} + (\frac{\sqrt{3}}{3}x + \frac{2\sqrt{6}}{3})^2 = 1$ \n    展开并化简：\n    $\frac{x^2}{9} + (\frac{3}{9}x^2 + 2 \cdot \frac{\sqrt{3}}{3}x \cdot \frac{2\sqrt{6}}{3} + \frac{4 \cdot 6}{9}) = 1$ \n    $\frac{x^2}{9} + \frac{x^2}{3} + \frac{4\sqrt{18}}{9}x + \frac{24}{9} = 1$ \n    $\frac{x^2}{9} + \frac{3x^2}{9} + \frac{4 \cdot 3\sqrt{2}}{9}x + \frac{8}{3} = 1$ \n    $\frac{4x^2}{9} + \frac{4\sqrt{2}}{3}x + \frac{8}{3} - 1 = 0$ \n    $\frac{4x^2}{9} + \frac{12\sqrt{2}}{9}x + \frac{5}{3} = 0$ \n    两边同乘以9，得到关于$x$的一元二次方程：\n    $4x^2 + 12\sqrt{2}x + 15 = 0$ \n\n4. **计算弦长（方法一：利用韦达定理和弦长公式）**：\n    设直线与椭圆的两个交点为$A(x_1, y_1)$和$B(x_2, y_2)$。\n    根据韦达定理，对于方程$4x^2 + 12\sqrt{2}x + 15 = 0$，有：\n    $x_1 + x_2 = -\frac{12\sqrt{2}}{4} = -3\sqrt{2}$ \n    $x_1 x_2 = \frac{15}{4}$ \n    弦长公式为$|AB| = \sqrt{(x_2 - x_1)^2 + (y_2 - y_1)^2}$。\n    由于直线斜率$k = \frac{y_2 - y_1}{x_2 - x_1} = \frac{\sqrt{3}}{3}$，所以$y_2 - y_1 = k(x_2 - x_1)$。\n    因此，$|AB| = \sqrt{(x_2 - x_1)^2 + k^2(x_2 - x_1)^2} = \sqrt{1 + k^2} |x_2 - x_1|$。\n    首先计算$(x_2 - x_1)^2 = (x_1 + x_2)^2 - 4x_1x_2$：\n    $(x_2 - x_1)^2 = (-3\sqrt{2})^2 - 4(\frac{15}{4}) = (9 \cdot 2) - 15 = 18 - 15 = 3$ \n    所以$|x_2 - x_1| = \sqrt{3}$。\n    接着计算$\sqrt{1 + k^2}$：\n    $\sqrt{1 + (\frac{\sqrt{3}}{3})^2} = \sqrt{1 + \frac{3}{9}} = \sqrt{1 + \frac{1}{3}} = \sqrt{\frac{4}{3}} = \frac{2}{\sqrt{3}} = \frac{2\sqrt{3}}{3}$ \n    最后，弦长$|AB| = \frac{2\sqrt{3}}{3} \cdot \sqrt{3} = \frac{2 \cdot 3}{3} = 2$。\n\n5. **计算弦长（方法二：利用焦点弦长公式）**：\n    对于过焦点的弦，弦长$|AB| = \frac{2a(1-e^2)}{1-e^2\cos^2\alpha}$，其中$e$是椭圆的离心率，$\alpha$是焦点弦所在的直线与$x$轴正方向的夹角。\n    离心率$e = \frac{c}{a} = \frac{2\sqrt{2}}{3}$。\n    直线的倾斜角是$\frac{\pi}{6}$，所以$\alpha = \frac{\pi}{6}$。\n    $\cos^2\alpha = \cos^2(\frac{\pi}{6}) = (\frac{\sqrt{3}}{2})^2 = \frac{3}{4}$。\n    $1-e^2 = 1 - (\frac{2\sqrt{2}}{3})^2 = 1 - \frac{8}{9} = \frac{1}{9}$。\n    将上述值代入焦点弦长公式：\n    $|AB| = \frac{2a(1-e^2)}{1-e^2\cos^2\alpha} = \frac{2 \cdot 3 \cdot (\frac{1}{9})}{1 - \frac{8}{9} \cdot \frac{3}{4}}$ \n    $= \frac{\frac{6}{9}}{1 - \frac{2}{3}} = \frac{\frac{2}{3}}{\frac{1}{3}} = 2$。\n    两种方法都得到弦长为$2$。\n\n最终答案是$2$。}


\end{document}
